\section{Rank coordinates and limiting measures}\label{sec:rank-limits}
We now suppose, for a contradiction, that there is a sequence of batch
systems for which
\[
 N/a\longrightarrow\infty,\qquad R\asymp(N/a)^{4/3}.
\]
Here each of the $R$ batches exchanges $a$ lost points for $a$ gained
points, and contributes the $a^2$ rightward edges between them.  Write
\[
 E_*=Ra^2,\qquad M=aN.
\]
Thus $E_*^3\asymp a^2N^4$ and $E_*/M\to\infty$.  The constants below may
depend on the two constants in the displayed comparison for $R$.
Our objective in this section is to assign compact labels to points so
that edge mass is controlled by vertex mass, and every positive portion
of edge mass sees crossings between different time windows.

Give batch $i$ the time $(i-1/2)/R$.  A \emph{window} is a dyadic
subinterval of $[0,1]$, with the half-open convention except at the right
endpoint $1$.  Its edges are those of its batches.  The consecutive
batches in a nonempty window have a slope span, from the state just
before its first batch to the state just after its last batch.  Disjoint
windows have slope spans with disjoint interiors.  Every edge slope is
strictly inside its batch span.  Each fixed window is nonempty eventually.
For a fixed time $t\in(0,1)$, choose batch $i=1+\lfloor Rt\rfloor$
and take the arithmetic mean of its two endpoint-state slopes.
A \emph{line of direction $t$} means a line parallel to this
configuration-dependent slope. The half-open convention fixes the
choice at cell boundaries. For fixed times and windows separated by
positive gaps, these directions have the same strict order as their
times for all sufficiently late configurations.

Let $S_J$ denote the number of distinct endpoints of the exchange edges
in $J$. These are exactly the points whose membership changes between
the endpoint states of at least one batch in $J$. A point that changes
membership within a batch and returns before its end contributes to
the sweep activity but need not contribute to $S_J$. Thus $S_J$ is
at most the activity of the full sweep interval covered by $J$.
The finite estimates proved above give
\begin{equation}\label{eq:finite-window-crossings}
 \operatorname{cr}(J)\le C a^2S_J^2.
\end{equation}
They also give convex and concave partitions of the full edge set into
at most $M$ chains; the edges of a window have such partitions into at
most $aS_J$ chains.

\subsection{A pair of rank coordinates parametrizes the plane}
For a nonempty window $J$, extend every chain in its convex partition to
a full convex piecewise linear graph.  On the left choose a ray slope
between the left endpoint of the window span and the first edge slope;
on the right choose one between the last edge slope and the right
endpoint of that span.  Add affine graphs until there are exactly $M$
graphs, denoted $g_{J,1},\ldots,g_{J,M}$.  All their slopes lie in a
closed interval $[\ell_J,u_J]$ strictly inside the window span.  The
finitely many choices of ray slopes and affine translates can avoid
all proper crossings between an edge in $J$ and an edge outside $J$.
At each such crossing exactly one of these $M$ graphs passes through
the crossing point: its actual edge graph.  This uses the genericity
assumption excluding a third determined segment through a proper
crossing.  For an empty window use arbitrary affine graphs; assertions
about any fixed finite collection of windows are made after those
windows have become nonempty.

Choose a continuous strictly positive probability density on $\mathbb R$
with distribution function $\Phi$.  For a sufficiently small
$\varepsilon_J>0$ put
\[
 r_J(x,y)=\frac1M\sum_{i=1}^M
 \Phi\!\left(\frac{y-g_{J,i}(x)}{\varepsilon_J}\right).
\]
This is the expected fraction of graphs below $(x,y)$ after a random
vertical translation.  Choose $\varepsilon_J$ so that, at every proper
crossing under consideration, the contribution of all graphs other
than the one through the point differs from their unsmoothed count by
at most $1$.  The choice is possible because there are finitely many
crossings and all these other vertical distances are nonzero.  Hence
at a crossing on the $i$th graph in height order,
\begin{equation}\label{eq:rank-slot-error}
 |r_J-i/M|\le 2/M.
\end{equation}
The choices can be made separately for all countably many windows.

For fixed $x$, the function $y\mapsto r_J(x,y)$ is continuous, strictly
increasing, and has limits $0$ and $1$.  Every level $r_J=v$, where
$0<v<1$, is therefore a full continuous graph $y=g_{J,v}(x)$.
For $x'>x$ the inequalities for the slopes of the original graphs give
\[
 \ell_J(x'-x)\le g_{J,v}(x')-g_{J,v}(x)
                    \le u_J(x'-x).
\]
For example, translating $(x,y)$ by $(h,u_Jh)$ can only increase $r_J$;
the upper inequality follows by applying this at a level point.  The
lower inequality is analogous.

Suppose $A$ precedes the disjoint window $B$.  Then $u_A<\ell_B$.
Consequently the difference of any $A$-level graph and any $B$-level
graph is strictly decreasing, with opposite infinite limits at the
two ends of the real line.  The graphs meet exactly once.  Their
intersection depends continuously on the two levels: bracket its
$x$-coordinate on either side by strict inequalities, which persist
under small changes of the levels, and then use continuity of a level
graph.  We have proved that
\begin{equation}\label{eq:rank-homeomorphism}
 (r_A,r_B):\mathbb R^2\longrightarrow(0,1)^2
 \quad\hbox{is a homeomorphism}.
\end{equation}
If $C$ is disjoint from both $A$ and $B$, expressing $r_C$ in these
coordinates gives a function monotone in each coordinate.  Indeed,
along an $A$-level graph, $r_B$ decreases strictly as $x$ increases;
along a $B$-level graph, $r_A$ increases strictly.  Compare the slope
range of $C$ with that of the level graph to determine the sign of
$r_C$.  The two signs, in the order $(r_A,r_B)$, are $(+,-)$ when
$C<A$, $(+,+)$ when $A<C<B$, and $(-,+)$ when $B<C$.

\subsection{Crossing density and variation}
For two disjoint windows $I,J$, sort their respective graph families
pointwise into $M$ ordered graphs.  Sorting preserves continuity and
the two bounds on chord slopes.  Each pair of sorted graphs, one from
each family, therefore meets at most once.  A proper crossing of an
$I$-edge and a $J$-edge determines a unique such pair of slots, and no
two crossings determine the same pair.  By
\eqref{eq:rank-slot-error}, the pair $(r_I,r_J)$ at this crossing is
within $2/M$ in each coordinate of the corresponding grid point
$(i/M,j/M)$.  After giving every crossing mass $M^{-2}$, every weak
limit of the pair-coordinate projection is thus at most planar
Lebesgue measure on $[0,1]^2$.  One way to see the last assertion is to
test a nonnegative continuous function: its integral is bounded by the
full $M\times M$ Riemann sum, up to its modulus of continuity at $4/M$.

There is also a uniform bound on the variation of a rank coordinate
along the edges:
\begin{equation}\label{eq:rank-total-variation}
 \sum_e\operatorname{Var}_e(r_J)\le 2M.
\end{equation}
To prove it, compare one vertically translated convex graph with one
chain of the concave partition of the full edge set.  Their height
difference is concave, so its positive set is an interval.  The
indicator of being above the translated graph has variation at most
$2$ along this chain.  Averaging over the translations and the $M$
convex graphs bounds the variation of $r_J$ along a concave chain by
$2$.  There are at most $M$ concave chains.  This proves the displayed
bound, including tangencies, because the superlevel-set argument
counts changes of sign rather than transverse intersections.

Enumerate the windows as $J_1,J_2,\ldots$ and give the compact cube
\[
 Z=[0,1]^{\mathbb N}
 \quad\hbox{the metric}\quad
 d(z,z')=\sum_{j\ge1}2^{-j}|z_{J_j}-z'_{J_j}|.
\]
Label every point $q$ of the plane by $z(q)=(r_J(q))_J$.
For an edge $e$, let $\operatorname{osc}_e z$ be the diameter of its
image in this metric.  By \eqref{eq:rank-total-variation} and summing
the nonnegative series,
\[
 \sum_e\operatorname{osc}_e z\le 2M.
\]
In particular, for each fixed $\rho>0$ the number of edges whose label
oscillation exceeds $\rho$ is at most $2M/\rho=o(E_*)$.

\subsection{The limiting measures}
Vertices, edges, and crossings require different normalizations.
We attach the rank label $z(q)$ to each vertex and crossing point;
an edge $e$ receives the label $(z(q_e),t_e)$, where $q_e$ is its left
endpoint and $t_e$ its batch time. The following table defines the
finite measures and names their weak limits:
\begin{center}
\begin{tabular}{lll}
Object & Weight of one finite object & Weak limit\\[3pt]
Vertex & $N^{-1}$ & $\mu$ on $Z$\\
Participating vertex in $J$ & $N^{-1}$ & $\mu_J$ on $Z$\\
Edge with its batch time & $E_*^{-1}$ & $\nu$ on $Z\times[0,1]$\\
Crossing between $I$ and $J$ & $M^{-2}$ & $\sigma_{IJ}$ on $Z$
\end{tabular}
\end{center}
Here $E_*=Ra^2$ and $M=aN$, and crossing measures are defined for
disjoint windows $I,J$. Each participating vertex is counted once
within a window. Thus, for example, the full vertex and edge measures
are $N^{-1}\sum_{q\in P}\delta_{z(q)}$ and
$E_*^{-1}\sum_e\delta_{(z(q_e),t_e)}$.
Compactness and bounded masses give a common weakly convergent
subsequence for this countable collection of measures. In particular
$\mu$ and $\nu$ are probability measures. Define the endpoint marginals
and the total participating-vertex masses by
\[
 s_J=\mu_J(Z),\qquad
 \lambda=\pi_{Z\,*}\nu,\qquad
 \lambda_J(F)=\nu(F\times J).
\]
The measures $\lambda$ and $\lambda_J$ count edge endpoints with edge
weights, rather than counting distinct vertices. In particular,
$\lambda_J$ is additive over disjoint time windows, whereas the
participating sets counted by $\mu_J$ may overlap.
We have $\mu_J\le\mu$ and $S_J/N\to s_J$.
The time marginal of $\nu$ is Lebesgue measure, since each batch has
the same number $a^2$ of edges and the batch times form a uniform grid.
In particular every fixed window boundary is $\nu$-null.  Restrictions
to windows therefore converge weakly, and $\lambda_J(Z)=|J|$.
Disintegration on the compact metrizable spaces~\cite[Example~2 and Theorem~3]{VianaRokhlin} supplies probability
measures $h_z$ on $[0,1]$, measurable in $z$, such that
\begin{equation}\label{eq:conditional-time}
 \lambda_J=h_z(J)\lambda.
\end{equation}
Equivalently, $\lambda_J(F)=\int_F h_z(J)\,d\lambda(z)$ for every
Borel $F\subset Z$. Thus $h_z$ is the conditional distribution of batch
time under edge mass, given the endpoint rank label.
On the $\lambda$-null set where a version has not been specified, choose
any probability measure.  Finally, the preceding grid argument gives
\begin{equation}\label{eq:crossing-projection}
 (z_I,z_J)_*\sigma_{IJ}\le\mathcal L^2\big|_{[0,1]^2}.
\end{equation}

The crossing projection bound concerns $\sigma_{IJ}$, whereas our
eventual contradiction must control the edge mass $\lambda$ using
the participating-vertex measures $\mu_J$. We next establish the
relations between these measures. For
any dyadic partition into windows of length $\delta$, the finite
activity-packing bound implies
\begin{equation}\label{eq:activity-tail}
 \#\{J: |J|=\delta,\ s_J>u\}\le C u^{-4/3}
 \qquad(u>0).
\end{equation}
Indeed a window with $S_J>uN$ has activity exceeding $uN$, and these
windows have disjoint step sets.  Apply the uniform finite estimate
and then take the limit for this finite collection of windows.

\begin{lemma}[Edge mass and vertex mass]\label{lem:edge-vertex-mass}
For every window $J$ and every Borel set $F\subset Z$,
\begin{equation}\label{eq:local-mass}
 \lambda_J(F)^3\le C s_J^3\mu_J(F).
\end{equation}
In particular $\lambda\ll\mu$.  The same estimate holds if additional
labels in a compact metrizable space are attached to vertices,
inherited by their edge endpoints, and a further weak limit is taken.
\end{lemma}
\begin{proof}
For a set $U$ of labels in a finite configuration, take the edges in
$J$ whose left endpoint belongs to $U$.  Write $e(U)$ for their number
and $m(U)$ for their number of distinct left endpoints.  The graph has
at most $S_J$ vertices, and all its edges have a designated endpoint
in a set of $m(U)$ vertices.  The designated-endpoint crossing
inequality, with its sparse-graph alternative, and
\eqref{eq:finite-window-crossings} give
\[
 e(U)^3\le C m(U)S_J\operatorname{cr}(J)+C S_J^3
       \le C a^2m(U)S_J^3+C S_J^3.
\]
After division by $E_*^3\asymp a^2N^4$, the last term tends to zero,
since it is at most $C/(a^2N)$.  Also $m(U)/N$ is at most the
participating-vertex measure of $U$.

For a compact $F$, use its open $\rho$-neighborhood as $U$ and bound
its vertex measure by that of the closed $2\rho$-neighborhood.
Weak convergence and the open-set and closed-set inequalities give
\[
 \lambda_J(F)^3\le C s_J^3\mu_J(F^{[2\rho]}).
\]
Let $\rho\downarrow0$.  Inner regularity of $\lambda_J$ then extends
the estimate from compact to all Borel sets.  Take $J=[0,1]$ to obtain
absolute continuity.  The finite argument makes no use of the nature
of the labels, proving the last assertion as well.
\end{proof}

\begin{lemma}[No atoms in conditional time]\label{lem:nonatomic-time}
For $\lambda$-almost every $z$, the probability measure $h_z$ is
nonatomic.
\end{lemma}
\begin{proof}
Fix $\beta>0$ and a dyadic partition of mesh $\delta$.  Put
$A_J=\{z:h_z(J)>\beta\}$.  At each $z$ at most $1/\beta$ of these
sets occur.  Thus $\sum_J\mu_J(A_J)\le1/\beta$.
For $M_0\ge1$, the windows with $s_J>M_0\delta^{3/4}$ contribute at
most $C M_0^{-4/3}$ to $\sum_J\lambda_J(A_J)$, by
\eqref{eq:activity-tail} and $\lambda_J(Z)=\delta$.
The remaining windows contribute, by
\eqref{eq:local-mass} and H\"older's inequality, at most
\[
 C M_0\delta^{3/4}\sum_J\mu_J(A_J)^{1/3}
 \le C M_0\beta^{-1/3}\delta^{1/12}.
\]
First let $\delta\downarrow0$, then $M_0\to\infty$.
If $h_z$ has an atom of mass greater than $\beta$, the member of every
partition containing that atom contributes at least its mass to
$\sum_Jh_z(J)\mathbf1_{A_J}(z)$.  Fatou's lemma therefore shows that
such $z$ form a $\lambda$-null set.  A countable sequence of positive
$\beta$ tending to zero proves the claim.
\end{proof}

\begin{lemma}[Edge mass sees crossings]\label{lem:edge-crossing-ac}
For every window $K$,
\begin{equation}\label{eq:crossing-ac}
 \lambda_K\ll\sum_{I,J\subset K,\ I\cap J=\varnothing}\sigma_{IJ},
\end{equation}
where the sum specifies null sets.  Thus, if a Borel set $F\subset Z$ is null for every
$\sigma_{IJ}$ with $I,J\subset K$ disjoint windows, then
$\lambda_K(F)=0$.
\end{lemma}
\begin{proof}
By inner regularity it suffices to consider a compact set $F$ that is
null for every crossing measure on the right of \eqref{eq:crossing-ac}.
Fix a dyadic resolution $\delta$ subdividing $K$, and discard windows with
$s_J>\delta^{2/3}$.  Their number is at most $C\delta^{-8/9}$, so the
limiting edge mass lost is at most $C\delta^{1/9}$.
In the remaining windows keep edges whose left-endpoint label lies in
the open $\rho$-neighborhood of $F$ and whose label oscillation is at
most $\rho$.  For fixed $\delta,\rho>0$, the oscillation condition
costs $o(E_*)$ edges.  If $e_{\delta,\rho}$ is the total number retained,
weak convergence consequently gives
\[
 \liminf \frac{e_{\delta,\rho}}E_*
 \ge \lambda_K(F)-C\delta^{1/9}.
\]
A crossing between retained edges from different windows has label
in the closed $2\rho$-neighborhood of $F$.  For this fixed finite
partition, its normalized limsup is therefore bounded by
\[
 \sum_{I\ne J}\sigma_{IJ}(F^{[2\rho]}),
\]
where each unordered pair is counted once.  This tends to zero as
$\rho\downarrow0$.  Crossings within one retained window are bounded
by \eqref{eq:finite-window-crossings}; after division by $M^2$ their
total limsup is at most
\[
 C\sum_{J:s_J\le\delta^{2/3}}s_J^2
 \le C\delta^{-1}\delta^{4/3}=C\delta^{1/3}.
\]
The ordinary crossing inequality, including its sparse alternative,
reads $e_{\delta,\rho}^3\le C N^2\operatorname{cr}+C N^3$.
Divide by $E_*^3\asymp a^2N^4$, take the configuration limit, and then
let $\rho\downarrow0$.  It follows that
\[
 \bigl(\lambda_K(F)-C\delta^{1/9}\bigr)_+^3
 \le C\delta^{1/3}.
\]
Finally let $\delta\downarrow0$.  This gives $\lambda_K(F)=0$, as
required.  Equivalently, $\lambda_K$ is absolutely continuous with
respect to any countable mixture of the crossing measures in the
statement whose coefficients are all positive and whose total mass is finite.
\end{proof}

The compact labels now retain three useful features: two separated
rank coordinates parametrize the plane; conditional edge activity
has no time atoms; and every positive portion of activity is detected
by a crossing measure with a bounded planar projection.  These are
the inputs to the local geometric argument.
